\section{Công cụ tự học và tổng kết}
\label{app:learning}

\subsection*{D-Tracker trong một trang}

\begin{keypoint}
\textbf{Đầu vào:} cửa sổ tensor $k\times l\times L_c$ (từ khóa $\times$ địa điểm $\times$ thời gian).
\textbf{Đầu ra:} dự báo $L_f$ bước và một mô hình có thể đọc được.
\textbf{Cách làm:} tách dữ liệu thành trend + seasonal + outlier; nén trend thành lõi nhỏ; cho lõi
tiến hóa bằng reaction--diffusion; dùng ALS để ước lượng; dùng MDL để chọn số chiều, outlier và
regime; cập nhật theo cửa sổ.
\end{keypoint}

\begin{enumerate}
  \item \textbf{Vì sao khó?} Dữ liệu đa chiều, có trend/mùa vụ/diffusion/outlier, các pattern lại
        đổi theo thời gian và stream gần như vô hạn.
  \item \textbf{Vì sao tách ba thành phần?} Để không nhầm mùa vụ với trend và không để một cú sốc
        làm hỏng toàn bộ hệ số.
  \item \textbf{Vì sao nén trend?} Để biến hàng chục nghìn quan sát thành vài dynamics ẩn và hai
        bảng hệ số từ khóa/địa điểm.
  \item \textbf{Vì sao reaction--diffusion?} Để biểu diễn cả tăng trưởng cục bộ lẫn lan truyền giữa
        các nhóm địa điểm bằng tham số có ý nghĩa.
  \item \textbf{Vì sao ALS?} Vì tối ưu đồng thời mọi khối là khó; cập nhật lần lượt dễ triển khai
        hơn và giữ được các ràng buộc.
  \item \textbf{Vì sao MDL?} Để chọn độ phức tạp có căn cứ thay vì thử tay, đồng thời quyết định
        outlier và regime.
  \item \textbf{Vì sao streaming được?} Mỗi bước chỉ dùng cửa sổ $L_c$ cố định, không quét lại toàn
        bộ lịch sử.
  \item \textbf{Kết quả?} Paper báo cáo MAE tốt nhất trên sáu dataset, thời gian gần như không tăng
        theo chiều dài stream, và các nhóm diffusion đọc được; không chứng minh nhân quả.
\end{enumerate}

\subsection*{Lộ trình đọc lại paper}

\begin{center}
\small
\begin{tabular}{@{}>{\raggedright\arraybackslash}p{0.27\linewidth}>{\raggedright\arraybackslash}p{0.31\linewidth}>{\raggedright\arraybackslash}p{0.34\linewidth}@{}}
  \toprule
  Đọc gì & Mục tiêu & Tự kiểm tra \\
  \midrule
  Abstract + §1 & Nắm bài toán, ba tính chất và Figure~1 & Kể được ba thách thức bằng lời \\
  §3.1 + paper Table~2 & Hiểu ký hiệu và Problem~1 & Viết input/output bằng ký hiệu \\
  §3.2--3.3 + Eq.~(1)--(4) & Hiểu decomposition và ODE & Giải thích từng số hạng \\
  §4.1 + Eq.~(5)--(9) & Hiểu ALS, LM, multiplicative update & Nêu thứ tự cập nhật \\
  §4.2--4.3 + Eq.~(10) & Hiểu MDL và ModelUpdate & Giải thích đánh đổi fit/complexity \\
  §5 + paper Table~3 + Figure~3--5 & Đọc kết quả và ablation & Phân biệt MAE/RMSE, claim và bằng chứng \\
  Appendix A.1--A.4 & Chi tiết khởi tạo và thuật toán & Chỉ ra các chỗ paper mơ hồ \\
  \bottomrule
\end{tabular}
\end{center}

Gợi ý thời lượng cho lần đọc đầu: 20 phút cho §1--3, 35 phút cho §4, 25 phút cho §5, 20 phút cho
phụ lục. Không cần hiểu hết các phép biến đổi đại số ngay lần đầu; hãy viết lại mỗi công thức thành
một câu tiếng Việt trước.

\subsection*{Bảng ký hiệu}

\begin{longtable}{@{}>{\raggedright\arraybackslash}p{0.23\linewidth}>{\raggedright\arraybackslash}p{0.35\linewidth}>{\raggedright\arraybackslash}p{0.35\linewidth}@{}}
  \caption{Ký hiệu, kích thước và nghĩa.}\label{tab:appendix-symbols}\\
  \toprule
  Ký hiệu & Kích thước & Nghĩa \\
  \midrule
  \endfirsthead
  \toprule
  Ký hiệu & Kích thước & Nghĩa \\
  \midrule
  \endhead
  $k,l$ & số nguyên & số từ khóa, số địa điểm \\
  $t_c$ & thời điểm & thời điểm hiện tại \\
  $L_c,L_f$ & số bước & độ dài cửa sổ, tầm dự báo \\
  $\mathcal{X}^{c}$ & $k\times l\times L_c$ & tensor quan sát hiện tại \\
  $d_k,d_l$ & số nguyên & số nhóm từ khóa, địa điểm \\
  $d_s$ & số nguyên & số pattern mùa vụ \\
  $\mathcal{W}^{(\mathrm{core})}$ & $d_k\times d_l\times L_c$ & lõi trend ẩn \\
  $\mathbf{W}^{(\mathrm{key})}$ & $d_k\times k$ & hệ số từ khóa, không âm \\
  $\mathbf{W}^{(\mathrm{loc})}$ & $d_l\times l$ & hệ số địa điểm, không âm \\
  $\mathbf{A}$ & $d_k\times d_l$ & tốc độ tăng trưởng $a_{ij}$ \\
  $\mathcal{D}$ & $d_k\times d_l\times d_l$ & độ mạnh diffusion $d_{ijj'}$ \\
  $\theta$ & --- & $\{\mathbf{A},\mathcal{D}\}$ \\
  $\theta_d,\theta_s$ & --- & tham số trend, tham số mùa vụ \\
  $\Theta$ & --- & một mô hình đầy đủ \\
  $\mathcal{F}$ & --- & tập các mô hình theo regime \\
  $\widehat{\mathcal{X}}_d,\widehat{\mathcal{X}}_s,\widehat{\mathcal{X}}_o$ & $k\times l\times L_c$ & trend, seasonal, outlier \\
  $\mathbf{S}^{(\mathrm{time})}$ & $d_s\times L_c$ & pattern mùa vụ theo thời gian \\
  $\mathbf{S}^{(\mathrm{key})},\mathbf{S}^{(\mathrm{loc})}$ & $d_s\times k$, $d_s\times l$ & hệ số mùa vụ \\
  \bottomrule
\end{longtable}

\subsection*{Thuật ngữ Việt--Anh}

\begin{longtable}{@{}>{\raggedright\arraybackslash}p{0.27\linewidth}>{\raggedright\arraybackslash}p{0.33\linewidth}>{\raggedright\arraybackslash}p{0.33\linewidth}@{}}
  \caption{Thuật ngữ dùng trong tài liệu.}\label{tab:glossary}\\
  \toprule
  Tiếng Việt & English & Ghi chú ngắn \\
  \midrule
  \endfirsthead
  \toprule
  Tiếng Việt & English & Ghi chú \\
  \midrule
  \endhead
  tensor & tensor & mảng nhiều chiều \\
  luồng tensor & tensor stream & chuỗi tensor theo thời gian \\
  cửa sổ & window & phần gần nhất dùng để ước lượng \\
  tầm dự báo & forecast horizon & số bước cần dự báo \\
  phân rã tensor & tensor factorization & biểu diễn tensor bằng các nhân tố nhỏ \\
  lõi & core & tensor nhỏ chứa dynamics ẩn \\
  nhân tố/hệ số & factor & mức liên hệ của đối tượng với nhóm \\
  động học ẩn & latent dynamics & chuỗi ẩn mô tả pattern \\
  phản ứng & reaction & tăng/giảm cục bộ \\
  khuếch tán & diffusion & lan truyền giữa các nhóm \\
  mùa vụ & seasonality & pattern lặp theo chu kỳ \\
  ngoại lai & outlier & giá trị bất thường, thưa \\
  chế độ & regime & giai đoạn có cấu trúc tương đối ổn định \\
  bình phương tối thiểu xen kẽ & alternating least squares (ALS) & tối ưu từng khối \\
  mô tả ngắn nhất & minimum description length (MDL) & chọn mô hình theo tổng số bit \\
  không âm & non-negative & $\ge 0$ \\
  dự báo theo luồng & stream forecasting & dự báo khi dữ liệu đến dần \\
  độ phức tạp & computational complexity & cách chi phí tăng theo kích thước \\
  \bottomrule
\end{longtable}

\subsection*{Câu hỏi tự kiểm tra phân tầng}

\paragraph{Tầng 1 --- Nhớ và hiểu.}
\begin{enumerate}
  \item Ba thành phần của D-Tracker là gì và mỗi thành phần xử lý loại pattern nào?
  \item Vì sao $\mathbf{W}^{(\mathrm{key})}$ và $\mathbf{W}^{(\mathrm{loc})}$ phải không âm?
  \item MDL gồm hai loại chi phí nào?
  \item Vì sao chỉ dùng cửa sổ gần nhất lại giúp xử lý stream?
\end{enumerate}

\paragraph{Tầng 2 --- Áp dụng và tính toán.}
\begin{enumerate}
  \item Với $d_k=3,d_l=2,L_c=104$, lõi có bao nhiêu phần tử? Hai bảng hệ số có bao nhiêu nếu
        $k=8,l=50$?
  \item Nếu $w_{j'}=3,w_j=1,d_{jj'}=0.4,a_j=0$, đạo hàm reaction--diffusion bằng bao nhiêu?
  \item Một mô hình tăng $d_s$ từ 2 lên 3: chi phí mô hình thay đổi theo hướng nào, và vì sao
        chất lượng fit có thể tăng?
\end{enumerate}

\paragraph{Tầng 3 --- Phản biện.}
\begin{enumerate}
  \item Nêu một cơ chế ngoài diffusion có thể tạo ra pattern trễ giữa hai quốc gia.
  \item Nếu chọn cửa sổ ngắn hơn chu kỳ mùa vụ, điều gì có thể xảy ra với thành phần seasonal?
  \item Vì sao ``MAE tốt nhất'' chưa đủ để kết luận mô hình tốt cho mọi quyết định?
  \item Cần thêm thí nghiệm gì để kiểm tra các nhóm địa điểm có ổn định không?
\end{enumerate}

\subsection*{Gợi ý trả lời}

\paragraph{Tầng 1.}
(1) Trend chứa thay đổi có cấu trúc và diffusion; seasonal chứa pattern lặp; outlier chứa bất thường
thưa. (2) Để phép chiếu không đảo dấu dynamics và giữ cách đọc ``nhóm này tăng/giảm''. (3) Chi phí
mô tả mô hình và chi phí mã hóa dữ liệu dưới mô hình. (4) Chi phí mỗi bước không phụ thuộc tổng
chiều dài lịch sử, và mô hình cũ có thể được thay khi pattern đổi.

\paragraph{Tầng 2.}
(1) Lõi $3\cdot2\cdot104=624$; hệ số $3\cdot8+2\cdot50=124$. (2)
$0+0.4(3-1)=0.8$. (3) Chi phí mô hình tăng vì thêm tham số; fit có thể tốt hơn vì biểu diễn được
nhiều pattern mùa vụ hơn, nhưng chỉ giữ nếu tổng MDL giảm.

\paragraph{Tầng 3.}
(1) Một chiến dịch toàn cầu, tin tức lan nhanh, hoặc cùng chịu ảnh hưởng của một mùa/ sự kiện.
(2) Seasonal có thể bị hấp thụ sai vào trend/outlier vì mô hình không thấy đủ một chu kỳ.
(3) MAE bỏ qua bất định, chi phí, khả năng diễn giải, drift và phân bố lỗi; quyết định rủi ro cần
thêm khoảng tin cậy và kiểm định. (4) Chạy nhiều khởi tạo/seed, so sánh nhóm qua các regime, và
kiểm tra độ ổn định của hệ số hoặc dùng dữ liệu held-out theo địa điểm.

\subsection*{Checklist trước khi nói ``đã hiểu paper''}

\begin{itemize}
  \item Tôi có thể vẽ lại pipeline từ tensor đến forecast mà không nhìn tài liệu.
  \item Tôi giải thích được Eq.~(1)--(4) bằng lời, kể cả kích thước các tensor.
  \item Tôi phân biệt được $\theta$, $\theta_d$, $\Theta$, $\mathcal{F}$.
  \item Tôi nêu được vai trò của LM, multiplicative update, ALS và MDL.
  \item Tôi đọc được paper Table~3 mà không nhầm đơn vị, và không biến ``57\%'' thành một kết luận chắc chắn.
  \item Tôi nói rõ được ba điều: kết quả paper, giả định của mô hình, và giới hạn của kết luận.
\end{itemize}
